\documentclass[12pt, a4paper]{article}

% --- Packages ---
% For xelatex/lualatex with Vietnamese support
\usepackage{fontspec}
\usepackage{polyglossia}
\setmainlanguage{vietnamese}
\setmainfont{DejaVu Serif}[
  BoldFont = DejaVu Serif Bold,
  ItalicFont = DejaVu Serif Italic,
  BoldItalicFont = DejaVu Serif Bold Italic
]
\setsansfont{DejaVu Sans}
\setmonofont{DejaVu Sans Mono}
\usepackage{geometry}
\usepackage{amsmath}
\usepackage{graphicx}
\usepackage{booktabs}
\usepackage{array}
\usepackage{hyperref}
\usepackage{setspace}
\usepackage{titlesec}
\usepackage{enumitem}
\usepackage{caption}
\usepackage{authblk}
\usepackage{abstract}
\usepackage{fancyhdr}
\usepackage{xcolor}
\usepackage{microtype}

\geometry{
  top=2.5cm, bottom=2.5cm,
  left=2.5cm, right=2.5cm
}

\hypersetup{
  colorlinks=true,
  linkcolor=blue,
  urlcolor=blue,
  citecolor=blue
}

\setlength{\headheight}{14pt}
\addtolength{\topmargin}{-2pt}

\pagestyle{fancy}
\fancyhf{}
\rhead{\small Z. Tang và cộng sự \quad \textit{Plant Phenomics} 7 (2025) 100025}
\rfoot{\thepage}

\titleformat{\section}{\normalfont\large\bfseries}{\thesection.}{0.5em}{}
\titleformat{\subsection}{\normalfont\normalsize\bfseries\itshape}{\thesubsection.}{0.5em}{}

\onehalfspacing

% ===================================================================
\begin{document}

% --- Tiêu đề ---
\begin{center}
  {\small \textit{Bài báo dữ liệu}}\\[6pt]
  {\LARGE \textbf{CVRP: Bộ dữ liệu hình ảnh lúa với chú thích chất lượng cao\\
  phục vụ phân đoạn ảnh và nghiên cứu kiểu hình thực vật}}\\[10pt]
  {\normalsize
    Zhiyan Tang$^{a,1}$, Jiandong Sun$^{a,1}$, Yunlu Tian$^{a,1}$,
    Jiexiong Xu$^{a}$, Weikun Zhao$^{a}$, Gang Jiang$^{a}$,
    Jiaqi Deng$^{a}$, Xiangchao Gan$^{a,b,*}$
  }\\[6pt]
  {\small
    $^{a}$Phòng thí nghiệm trọng điểm quốc gia về Di truyền, Tài nguyên và Cải tiến cây trồng,
    Đại học Nông nghiệp Nam Kinh, Nam Kinh 210095, Trung Quốc\\
    $^{b}$Phòng thí nghiệm Nhân giống sinh học Trung Sơn, Nam Kinh 210095, Trung Quốc\\[4pt]
    $^{1}$Các tác giả đóng góp ngang nhau.\\
    $^{*}$Tác giả liên hệ: \href{mailto:gan@njau.edu.cn}{gan@njau.edu.cn}
  }
\end{center}

\vspace{0.5cm}
\noindent\rule{\linewidth}{0.4pt}

% --- Thông tin bài báo & Tóm tắt ---
\noindent
\begin{minipage}[t]{0.28\linewidth}
  \small
  \textbf{Từ khóa:}\\
  Chú thích ảnh\\
  Bông lúa\\
  Phân đoạn ảnh\\
  Học sâu\\
  Kiểu hình thực vật
\end{minipage}%
\hfill
\begin{minipage}[t]{0.69\linewidth}
  \small
  \textbf{Tóm tắt}\\
  Các mô hình học máy phục vụ phân tích ảnh cây trồng và nghiên cứu kiểu hình có tầm quan
  trọng cao đối với nông nghiệp chính xác và công tác chọn giống, đồng thời là chủ đề được
  nghiên cứu chuyên sâu. Tuy nhiên, sự thiếu hụt các bộ dữ liệu hình ảnh chất lượng cao có
  chú thích chi tiết và được công bố công khai đã cản trở nghiêm trọng sự phát triển của các
  mô hình này. Trong công trình này, chúng tôi giới thiệu một bộ dữ liệu hình ảnh cây lúa đa
  giống và đa góc nhìn (CVRP) được tạo ra từ 231 giống địa phương và 50 giống hiện đại được
  trồng theo mật độ dày trên ruộng lúa nước. Bộ dữ liệu bao gồm ảnh chụp cây lúa trong môi
  trường tự nhiên ngoài đồng, cũng như ảnh trong phòng tập trung vào bông lúa, cho phép điều
  tra chi tiết các đặc điểm đặc trưng của từng giống. Quy trình chú thích bán tự động sử dụng
  mô hình học sâu đã được thiết kế để chú thích, sau đó được kiểm tra thủ công chặt chẽ.
  Chúng tôi chứng minh tính hữu ích của CVRP bằng cách đánh giá hiệu suất của bốn mô hình
  phân đoạn ngữ nghĩa tiên tiến nhất (SOTA). Chúng tôi cũng thực hiện tái tạo 3D cây trồng
  với phân đoạn bộ phận thông qua ảnh và chú thích. Cơ sở dữ liệu không chỉ hỗ trợ nhận
  diện và phân đoạn bông lúa theo mục đích chung mà còn cung cấp tài nguyên quý giá cho các
  nhiệm vụ khó khăn như nhận diện tự động giống lúa, đếm bông và hạt lúa, và tái tạo 3D cây
  trồng. Cơ sở dữ liệu và mô hình chú thích ảnh có tại
  \href{https://bic.njau.edu.cn/CVRP.html}{https://bic.njau.edu.cn/CVRP.html}.
\end{minipage}

\noindent\rule{\linewidth}{0.4pt}
\vspace{0.3cm}

% ===================================================================
\section{Giới thiệu}

Là một loại lương thực chủ yếu, lúa là nguồn dinh dưỡng quan trọng cho hơn nửa dân số thế
giới. Đặc điểm kiểu hình của bông lúa là một chỉ số quan trọng phản ánh sự sinh trưởng và
năng suất~\cite{chandra2020}, khiến nó trở nên hữu ích cao trong canh tác và chọn giống lúa.
Tuy nhiên, các phương pháp thu thập dữ liệu truyền thống vẫn phụ thuộc nhiều vào quan sát và
đo đạc thủ công. Kiểu hình thực vật học, tận dụng học máy và thị giác máy tính, đang được
ứng dụng ngày càng nhiều và đang trở thành động lực thúc đẩy nghiên cứu nông nghiệp chính
xác và chọn giống.

Các nhà nghiên cứu thường trích xuất các đặc điểm kiểu hình hoặc thống kê bằng cách áp dụng
học sâu vào ảnh, giúp giảm đáng kể chi phí thời gian và lao động~\cite{xiong2017}.

Trong những năm gần đây, nhiều phương pháp tiên tiến đã được áp dụng để nâng cao độ chính
xác và hiệu quả trong nghiên cứu kiểu hình lúa, với trọng tâm là các kỹ thuật phân đoạn
bông lúa. Chẳng hạn, xác định giai đoạn trỗ bông của lúa không chỉ quan trọng để cải thiện
tỷ lệ đậu hạt và đảm bảo năng suất cao mà còn hữu ích cho các mô hình sinh trưởng lúa. Bai
và cộng sự~\cite{bai2018} đã giải quyết thách thức tự động nhận diện giai đoạn trỗ bông bằng
cách chuyển đổi vấn đề thành phát hiện từng bông lúa riêng lẻ. Để mô tả chính xác và toàn
diện các đặc điểm của bông lúa trong ảnh, các nhà nghiên cứu đã sử dụng cả đặc trưng thị
giác truyền thống như màu sắc và gradient, lẫn đặc trưng học được thông qua mạng nơ-ron tích
chập (CNN). Hayat và cộng sự~\cite{hayat2020} đề xuất phương pháp học Bayesian không giám sát
để phân đoạn bông lúa trong ảnh UAV. Mô hình hỗn hợp Gaussian và phương pháp Monte Carlo theo
chuỗi Markov (MCMC) được phát triển để phát hiện pixel bông lúa mà không cần dữ liệu huấn
luyện có nhãn. Xiao và cộng sự~\cite{xiao2023} giới thiệu phương pháp kết hợp đặc trưng
skip-connection có trọng số (WSFF) và tích hợp vào mô hình UNet để tạo ra weighted
skip-connection UNet (WSUNet), nhằm cải thiện phân đoạn bông lúa từ UAV mà không tăng chi
phí tính toán.

Hầu hết các phương pháp nhận diện và phân đoạn bông lúa đều sử dụng các mô hình dựa trên
học sâu, đòi hỏi nhiều ảnh có chú thích chất lượng cao để huấn luyện. Bông lúa trong ảnh có
sự khác biệt đáng kể về màu sắc, kết cấu, kích thước và hình dạng do các yếu tố môi trường
đồng ruộng cực kỳ phức tạp, bao gồm phản chiếu ánh sáng từ nước, mất cân bằng ánh sáng do
thời tiết và nền ảnh lộn xộn~\cite{xiong2017}. Tệ hơn nữa, có nhiều thiết bị kiểu hình khác
nhau và dữ liệu ảnh thu thập thường bị giới hạn cho các mục đích cụ thể. Trong Zhou và cộng
sự~\cite{zhou2023}, thiết bị thu thập ảnh được gắn trên giá cố định và máy ảnh chụp thẳng
đứng xuống ở góc 90 độ. Wang và cộng sự~\cite{wang2021} sử dụng UAV để chụp ảnh ruộng lúa,
nhưng độ rõ nét của các ảnh này thấp. Ảnh thông lượng cao bị giới hạn ở góc nhìn tán lá từ
góc thẳng đứng hoặc chếch. Mặc dù chụp ảnh từ góc chếch có thể thu được nhiều bông và đặc
trưng phong phú hơn so với góc thẳng đứng, nhưng không thể chụp được tất cả bông của cây
trồng dày từ một góc nhìn cố định do bị che khuất. Hơn nữa, sự biến đổi đáng kể của điều
kiện đồng ruộng như ánh sáng, che khuất và sự hiện diện của cỏ dại ảnh hưởng đến chất lượng
ảnh thu được và độ chính xác của phân tích kiểu hình tiếp theo.

Sự thiếu hụt các bộ dữ liệu hình ảnh chất lượng cao có chú thích chi tiết và được công bố
công khai đã cản trọng nghiêm trọng việc áp dụng các thuật toán thị giác máy tính SOTA trong
kiểu hình thực vật và phát triển phương pháp kiểu hình thực vật. Trong những năm gần đây,
nhiều nghiên cứu đã thu thập bộ dữ liệu ảnh để nghiên cứu các yếu tố nông học ảnh hưởng đến
năng suất cây trồng, nhưng khi sử dụng để huấn luyện mô hình học sâu, chú thích là bắt buộc
và thường tốn kém. Bộ dữ liệu phát hiện đầu lúa mì toàn cầu (GWHD)~\cite{david2020,
david2021} được xây dựng và mở rộng như một bộ dữ liệu phát hiện đầu lúa mì đa dạng và chất
lượng cao. Bộ dữ liệu MinneApple~\cite{hani2020} được tạo ra như một bộ dữ liệu chuẩn để
phát hiện và phân đoạn táo. Wang và cộng sự~\cite{wang2021} xây dựng bộ dữ liệu ảnh lúa
nước, giới thiệu quy trình huấn luyện mô hình học sâu bán giám sát để hỗ trợ chú thích và
tinh chỉnh bộ dữ liệu. Tuy nhiên, trong nghiên cứu này chỉ có hai giống lúa được thu thập từ
ruộng thực nghiệm và sự đa dạng về diện mạo bông lúa giữa các giống khác nhau bị bỏ qua.
Vẫn còn thiếu các bộ dữ liệu công khai toàn diện để huấn luyện và đánh giá mô hình. Do đó,
một số phương pháp học sâu phân đoạn bông lúa chỉ được xác nhận trên bộ dữ liệu hạn chế.
Chẳng hạn, Xiong và cộng sự~\cite{xiong2017} chỉ kiểm tra 48 mẫu với độ chính xác phân đoạn
bông khoảng 80\%.

Trong bài báo này, chúng tôi xây dựng bộ dữ liệu hình ảnh cây lúa đa giống và đa góc nhìn
(CVRP). Vì đặc điểm kiểu hình của bông lúa chịu ảnh hưởng của thông tin di truyền và điều
kiện môi trường, ảnh của 231 giống địa phương và 50 giống hiện đại được chụp độc lập trên
ruộng lúa nước. CVRP bao gồm 2.303 ảnh đồng ruộng với mặt nạ chú thích tương ứng và 123
ảnh bông lúa trong phòng. Ban đầu chúng tôi chú thích một phần ảnh theo cách thủ công và
huấn luyện mô hình phân đoạn ngữ nghĩa. Sử dụng mô hình đã huấn luyện, chúng tôi gán nhãn
bông lúa trên các ảnh chưa có nhãn còn lại rồi kiểm tra thủ công các chú thích đó. Quy
trình này không chỉ nâng cao hiệu quả tạo ra CVRP mà còn cải thiện đáng kể chất lượng của
nó, biến CVRP trở thành tài nguyên quý giá cho nghiên cứu kiểu hình lúa trong tương lai.
Hơn nữa, chúng tôi đánh giá hiệu suất của các mô hình phân đoạn ngữ nghĩa SOTA trên bộ dữ
liệu. Bộ dữ liệu này rất có giá trị để huấn luyện các mô hình trí tuệ nhân tạo (AI) ước
tính các đặc điểm lúa thông thường, bao gồm mật độ bông và đặc điểm bông. Hơn nữa, vì ảnh
của mỗi giống lúa được chụp và chú thích độc lập, theo hiểu biết tốt nhất của chúng tôi, đây
là bộ sưu tập lớn đầu tiên giúp huấn luyện các mô hình AI cho các nhiệm vụ độ phân giải cao
hơn như nhận diện giống lúa và đếm hạt theo giống.

% ===================================================================
\section{Vật liệu và phương pháp}

\subsection{Thu thập dữ liệu}

Chúng tôi chụp ảnh cây lúa trên ruộng lúa nước ở Nam Kinh, Trung Quốc, và đưa các bông lúa
riêng lẻ về phòng thí nghiệm tại Đại học Nông nghiệp Nam Kinh để chụp ảnh trong phòng. Các
mẫu của chúng tôi bao gồm 231 giống địa phương và 50 giống hiện đại từ ruộng thực nghiệm.
Tám giống địa phương với đặc điểm bông đặc trưng được thể hiện trong Hình 1. Việc thu thập
dữ liệu được thực hiện chủ yếu trong giai đoạn chín. Do trồng dày, có sự che khuất đáng kể
giữa các cây lúa. Để thu được các đặc điểm đặc trưng theo giống của bông lúa và ghi lại đặc
điểm chi tiết của bông mỗi giống, chúng tôi sử dụng bốn tấm nhựa đặc trưng để tách các cây
xung quanh trong khi chụp ảnh. Các tấm được vẽ với những hoa văn thị giác phong phú và đặc
biệt để ngay cả mắt thường cũng có thể nhận ra góc nhìn của ảnh trực tiếp. Đối với một số
cây, các ảnh chụp đủ các góc nhìn đại diện và cho phép tái tạo toàn bộ cây trong không gian
3D. Chúng tôi sử dụng kỹ thuật chụp theo vòng xoắn ốc để đảm bảo bao phủ nhiều góc nhìn.
Kỹ thuật này cho phép chụp được nhiều đặc điểm bông đa dạng như màu sắc, kết cấu và kích
thước. Để đảm bảo khả năng ứng dụng rộng rãi hơn, chúng tôi cũng chụp nhiều ảnh bông lúa
trên ruộng lúa dày mà không có nền nhân tạo (sau đây gọi là cảnh đồng ruộng). Một trăm ảnh
đồng ruộng của các cây ở trạng thái sinh trưởng tự nhiên được thu thập để chú thích.

Ngoài ra, chúng tôi chụp ảnh trong phòng của từng bông lúa riêng lẻ với thông tin cấu trúc
chi tiết theo hai tư thế. Một bộ ảnh chụp bông trong trạng thái tự nhiên với đế tự chế để
giữ nguyên tư thế và cấu trúc gốc. Bộ còn lại cho thấy bông ở trạng thái được mở ra bằng
tay, với các nhánh được sắp xếp để lộ ra các hạt ẩn và cấu trúc bên trong. Để giảm nhiễu từ
nền lộn xộn, chúng tôi sử dụng nền đen sậm trong quá trình chụp ảnh để làm nổi bật các hạt
của bông và cải thiện độ chính xác của phân tích ảnh tiếp theo. Mỗi ảnh được chụp trong điều
kiện ánh sáng được kiểm soát để đảm bảo tính nhất quán và giảm bóng. Những ảnh trong phòng
này làm phong phú thêm CVRP với các ảnh cận cảnh độ phân giải cao không có nhiễu môi trường.

\subsection{Chuẩn hóa ảnh}

Một nhiệm vụ quan trọng trong việc tạo bộ dữ liệu bông là chuẩn hóa các ảnh đa dạng. Để
đảm bảo chụp toàn diện các chi tiết của từng cá thể, chúng tôi trước tiên chụp video của
các cây rồi trích xuất ảnh chất lượng cao. Tám ảnh được chọn cho mỗi cây, ngoại trừ một số
giống địa phương trong điều kiện sinh trưởng kém. Chúng tôi cố gắng bao gồm các góc nhìn
thẳng đứng, chếch, ngang và cận cảnh cho mỗi giống. Các ví dụ về ảnh với các góc nhìn khác
nhau được trình bày trong Hình 3. Góc nhìn thẳng đứng cung cấp góc nhìn từ trên xuống, làm
nổi bật sự sắp xếp và mật độ của bông, trong khi góc nhìn chếch và ngang cung cấp thông tin
về đặc điểm cấu trúc và chi tiết kết cấu của bông và cây. Góc nhìn cận cảnh cung cấp thông
tin chi tiết hơn về một bông lúa riêng lẻ. Ngoài ra, chúng tôi đảm bảo các ảnh được chọn có
chất lượng cao bằng cách loại bỏ những ảnh bị mờ. Ảnh mờ có thể do máy ảnh di chuyển, lấy
nét kém hoặc gió làm rung, và chúng có thể ảnh hưởng đến độ chính xác của chú thích và hiệu
suất của các mô hình phân đoạn. Các ảnh đa giống được chuẩn hóa về độ phân giải $1080 \times
1920$. Đối với ảnh bông lúa riêng lẻ trong phòng, độ phân giải được đặt ở mức $1500 \times
2007$.

\subsection{Chú thích ảnh}

Chú thích bộ dữ liệu là một nhiệm vụ tốn thời gian, đặc biệt với dữ liệu bông lúa có cấu
trúc phức tạp và che khuất không đều. Để tối ưu hóa quy trình chú thích, chúng tôi kết hợp
các phương pháp học sâu với các phương pháp thủ công và chia quy trình thành hai giai đoạn:
chú thích thủ công và dự đoán dựa trên mô hình. Quy trình chú thích đầy đủ cho dữ liệu đồng
ruộng được trình bày trong Hình 5.

Trong giai đoạn đầu, chúng tôi chú thích thủ công 807 ảnh đồng ruộng thông qua công cụ chú
thích tương tác LabelMe~\cite{russell2008}. Hình đa giác do công cụ này cung cấp đặc biệt
hiệu quả trong việc mô tả các đường viền phức tạp và tinh tế của bông lúa, đảm bảo chú thích
chính xác. Quá trình này tốn nhiều thời gian và thực tế mất tổng cộng 352 giờ. Quá trình chú
thích thủ công chi tiết này tạo nền tảng để huấn luyện mô hình phân đoạn ngữ nghĩa mạnh mẽ.

Trong giai đoạn thứ hai, một mạng nơ-ron sâu được huấn luyện trên dữ liệu đã được chú thích
thủ công. Chúng tôi sử dụng Mask2Former, vốn sử dụng chú ý che mặt nạ và module trích xuất
đặc trưng độ phân giải cao để tập trung vào các vùng mục tiêu, cho phép phát hiện và phân
đoạn chính xác bằng cách nắm bắt các chi tiết tinh tế. Backbone của mô hình là Swin
Transformer tiên tiến. Bộ tối ưu hóa là AdamW, tăng cường khả năng tổng quát hóa và hiệu
suất bằng cách tách rời phân rã trọng số khỏi cập nhật gradient. Trong quá trình huấn luyện,
kích thước batch được đặt là 2, tốc độ học được đặt là 0,0001. Sau đó chúng tôi sử dụng mạng
nơ-ron đã huấn luyện để nhận diện và phân đoạn bông lúa trên các ảnh chưa được gán nhãn còn
lại. Các chú thích được dự đoán sau đó được kiểm tra thủ công, và trong thực tế chỉ cần điều
chỉnh nhỏ. Thời gian xử lý trung bình cho mỗi ảnh khoảng 1 phút.

Bằng cách tích hợp học sâu vào quy trình chú thích, chúng tôi giảm đáng kể cả thời gian lẫn
chi phí lao động. Phương pháp lai ghép này không chỉ tăng tốc quy trình chú thích mà còn đảm
bảo chất lượng chú thích cao. Bộ dữ liệu kết quả với các bông lúa được gán nhãn chính xác
cung cấp dữ liệu chất lượng cao cho việc phát triển và đánh giá mô hình phân đoạn bông lúa.
Thống kê chi tiết của bộ dữ liệu CVRP được trình bày trong Bảng 1.

% --- Table 1 ---
\begin{table}[ht]
  \centering
  \caption{Thống kê chi tiết của bộ dữ liệu CVRP.}
  \label{tab:dataset_stats}
  \begin{tabular}{lccccc}
    \toprule
    \textbf{Danh mục} & \multicolumn{2}{c}{\textbf{Dữ liệu đồng ruộng}} & \textbf{Dữ liệu trong phòng} & \textbf{Tổng cộng} \\
    \cmidrule(lr){2-3}
    & Đa giống & Cảnh đồng ruộng & & \\
    \midrule
    Số ảnh       & 2203 & 100 & 123 & 2426 \\
    Số chú thích & 2203 & 100 & --  & 2303 \\
    \bottomrule
  \end{tabular}
\end{table}

% ===================================================================
\section{Ứng dụng trong nhiệm vụ phân đoạn}

\subsection{Khung triển khai và môi trường hoạt động}

Để phân đoạn ảnh cây lúa, chúng tôi triển khai một số mô hình phân đoạn ngữ nghĩa SOTA thông
qua framework MMSegmentation v1.2.2~\cite{mmseg2020}, một bộ công cụ phân đoạn ngữ nghĩa mã
nguồn mở được xây dựng trên PyTorch. Bộ công cụ này cung cấp một hộp công cụ chuẩn thống
nhất cho nhiều phương pháp phân đoạn ngữ nghĩa khác nhau. Bộ công cụ này mô-đun hóa framework
phân đoạn ngữ nghĩa thành các thành phần, cho phép người dùng dễ dàng xây dựng framework tùy
chỉnh bằng cách kết hợp các mô-đun khác nhau. Trong framework MMsegmentation~\cite{mmseg2020},
MMCV (Nền tảng Thị giác Máy tính OpenMMLab) được sử dụng để xử lý dữ liệu và huấn luyện mô
hình. Cấu hình của chúng tôi bao gồm GPU A4000, CUDA v11.3 và PyTorch v1.10.0. Chúng tôi
kiểm tra Mô hình Phân đoạn Bất kỳ (SAM)~\cite{kirillov2023} và Mô hình Phân đoạn Bất kỳ 2
(SAM2)~\cite{ravi2024} để phân đoạn hạt từ ảnh bông lúa. Hơn nữa, chúng tôi sử dụng framework
tái tạo mục tiêu 3D và phân đoạn đa cấp một điểm dừng (OSTRA)~\cite{xu2023} để kiểm tra tái
tạo mục tiêu cây đơn thông qua bộ dữ liệu CVRP.

\subsection{Phân đoạn ngữ nghĩa ảnh cây lúa}

Để đánh giá CVRP trong phân đoạn ngữ nghĩa, chúng tôi triển khai bốn mô hình hiệu suất cao:
DeepLab Phiên bản 3 Plus (DeepLabV3+)~\cite{chen2018}, masked-attention mask transformer
(Mask2Former)~\cite{cheng2022}, phân đoạn ngữ nghĩa với transformer (SegFormer)~\cite{xie2021},
và mạng dựa trên kernel (K-Net)~\cite{zhang2021}. Mỗi mô hình này cung cấp những điểm mạnh
độc đáo cho việc đánh giá của chúng tôi.

\textbf{DeepLabV3+} sử dụng module gộp kim tự tháp không gian atrous (ASPP) để nắm bắt ngữ
cảnh đa tỷ lệ dựa trên tích chập atrous, không trích xuất đặc trưng theo từng pixel. Thay
vào đó, nó tăng trường tiếp nhận, cho phép mỗi đầu ra tích chập nắm bắt phạm vi thông tin
rộng hơn. Module ASPP hoạt động như một bộ tăng cường đặc trưng bằng cách sử dụng nhiều
kernel tích chập với các tốc độ giãn khác nhau.

\textbf{Mask2Former}~\cite{cheng2022} được xây dựng trên một meta-framework đơn giản với bộ
giải mã transformer mới. Thành phần chính của nó là chú ý che mặt nạ, trích xuất đặc trưng
cục bộ bằng cách giới hạn cross-attention trong vùng mặt nạ được dự đoán. Mask2Former phát
triển module trích xuất đặc trưng độ phân giải cao để trích xuất đặc trưng độ phân giải cao
từ bản đồ đặc trưng độ phân giải thấp nhằm nắm bắt các chi tiết của vùng mục tiêu.

\textbf{SegFormer}~\cite{xie2021} kết hợp transformer với bộ giải mã perceptron đa lớp nhẹ;
bộ mã hóa của nó không cần nội suy mã hóa vị trí khi xử lý ảnh với các độ phân giải huấn
luyện khác nhau, cho phép dễ dàng thích ứng với các độ phân giải kiểm tra tùy ý mà không ảnh
hưởng đến hiệu suất. Cấu trúc phân cấp có thể tạo ra đặc trưng chi tiết ở độ phân giải cao
và đặc trưng thô ở độ phân giải thấp.

\textbf{K-Net}~\cite{zhang2021} là một kiến trúc sáng tạo cho phân đoạn ảnh giới thiệu tạo
kernel động để nắm bắt đặc trưng của các đối tượng khác nhau. Bằng cách sử dụng trích xuất
đặc trưng đa tỷ lệ, K-Net tích hợp đặc trưng từ nhiều cấp độ khác nhau, nâng cao hiệu suất
cho các đối tượng đa tỷ lệ trong các cảnh phức tạp.

Trong framework MMsegmentation~\cite{mmseg2020}, các backbone sử dụng khác nhau giữa các
mạng: DeepLabV3+ được xây dựng trên mạng nơ-ron tàn dư (ResNet), SegFormer được xây dựng
trên Mix Vision Transformer, và cả Mask2Former lẫn K-Net đều được xây dựng trên Swin
Transformer.

Chúng tôi áp dụng các mô hình phân đoạn trên vào CVRP, sử dụng 727 ảnh làm tập huấn luyện
và 80 ảnh làm tập xác nhận, để đánh giá hiệu suất theo một số chuẩn mực tiêu chuẩn. Các
metric tiêu chuẩn bao gồm precision, recall, F score và intersection over union (IoU). Các
metric này dựa trên bốn đại lượng cơ bản: true positives (TP), true negatives (TN), false
positives (FP) và false negatives (FN). Trong công trình của chúng tôi, TP đại diện cho các
pixel bông được phân loại đúng là bông, TN đại diện cho các pixel nền được phân loại đúng là
nền, FP đại diện cho các pixel nền bị phân loại sai là bông, và FN đại diện cho các pixel
bông bị phân loại sai là nền.

Precision tính tỷ lệ dự đoán đúng dương trên tổng số dự đoán dương của mô hình, phản ánh độ
chính xác của phân đoạn bông:
\begin{equation}
  \text{Precision} = \frac{TP}{TP + FP}
\end{equation}

Recall đo tỷ lệ dự đoán đúng dương trên tổng số dương thực tế. Trong công trình của chúng
tôi, nó chỉ ra khả năng của mô hình trong việc nắm bắt tất cả bông liên quan:
\begin{equation}
  \text{Recall} = \frac{TP}{TP + FN}
\end{equation}

F score là trung bình điều hòa của precision và recall, cung cấp một thước đo cân bằng xem
xét cả false positives và false negatives:
\begin{equation}
  \text{F score} = 2 \cdot \frac{\text{Precision} \cdot \text{Recall}}{\text{Precision} + \text{Recall}}
\end{equation}

IoU đặc biệt quan trọng trong việc đánh giá các mô hình phân đoạn vì nó định lượng sự chồng
lấp giữa phân đoạn được dự đoán và ground truth. Giá trị IoU cao hơn biểu thị hiệu suất mô
hình tốt hơn:
\begin{equation}
  \text{IoU} = \frac{TP}{TP + FP + FN}
\end{equation}

Các metric này rất quan trọng để đo lường hiệu quả của các mô hình học máy trong việc nhận
diện và phân đoạn các đối tượng trong ảnh. Vì chúng tôi đánh giá hiệu suất phân đoạn của mô
hình trên CVRP, chúng tôi chỉ tập trung vào việc phân biệt bông lúa với các danh mục khác.
Vì nền chiếm tỷ lệ lớn trong ảnh, chúng tôi đánh giá các metric hiệu suất riêng cho danh mục
bông mà không lấy trung bình với danh mục nền.

Kết quả của chúng tôi được trình bày trong Bảng 2. Hiệu suất của bốn mô hình phân đoạn tương
đối có thể so sánh được với nhau. Recall nhìn chung trên 85\%, cho thấy tất cả các mô hình
đều có thể nhận diện chính xác phần lớn bông lúa. IoU cho thấy Mask2Former có hiệu suất tốt
nhất trong phân đoạn bông.

% --- Table 2 ---
\begin{table}[ht]
  \centering
  \caption{Các metric hiệu suất của bốn mô hình phân đoạn SOTA trên dữ liệu đa giống.}
  \label{tab:seg_results}
  \begin{tabular}{lcccc}
    \toprule
    \textbf{Mô hình phân đoạn} & \textbf{Precision $\uparrow$} & \textbf{Recall $\uparrow$} & \textbf{F score $\uparrow$} & \textbf{IoU $\uparrow$} \\
    \midrule
    DeepLabV3+ & 79.83 & 86.55 & 83.05 & 71.02 \\
    Mask2Former & 81.12 & \textbf{88.49} & \textbf{84.64} & \textbf{73.38} \\
    SegFormer   & 78.76 & 88.11 & 83.17 & 71.19 \\
    K-Net       & \textbf{82.37} & 85.85 & 84.08 & 72.53 \\
    \bottomrule
  \end{tabular}
\end{table}

Chúng tôi cũng đánh giá các mô hình đã huấn luyện trên ảnh cảnh đồng ruộng để đánh giá khả
năng tổng quát hóa. Kết quả được trình bày trong Bảng 3. Kết quả precision và recall cho thấy
mô hình đã huấn luyện Mask2Former nhận diện hiệu quả phần lớn bông lúa với mức độ chính xác
thỏa đáng. Đáng chú ý, một số tình huống cụ thể gây ra thách thức cho mô hình của chúng tôi.
Lá lúa và cỏ dại đôi khi bị nhầm lẫn là bông lúa ở cả hai giai đoạn sinh trưởng. Màu sắc
tương tự của lá lúa, cỏ dại và bông lúa có thể tạo ra các tín hiệu thị giác gây nhầm lẫn
cho phân đoạn, dẫn đến bỏ sót nhãn bông hoặc phân loại sai cỏ dại là bông. Hơn nữa, phân
đoạn các bông lúa nhỏ chiếm phần nhỏ trong ảnh rất khó khăn.

% --- Table 3 ---
\begin{table}[ht]
  \centering
  \caption{Các metric hiệu suất của bốn mô hình phân đoạn SOTA trên dữ liệu cảnh đồng ruộng.}
  \label{tab:field_results}
  \begin{tabular}{lcccc}
    \toprule
    \textbf{Mô hình phân đoạn} & \textbf{Precision $\uparrow$} & \textbf{Recall $\uparrow$} & \textbf{F score $\uparrow$} & \textbf{IoU $\uparrow$} \\
    \midrule
    DeepLabV3+ & 82.73 & 74.18 & 76.89 & 62.82 \\
    Mask2Former & \textbf{89.39} & \textbf{84.60} & \textbf{86.49} & \textbf{76.56} \\
    SegFormer   & 79.28 & 84.08 & 80.91 & 68.24 \\
    K-Net       & 85.52 & 77.50 & 80.13 & 67.34 \\
    \bottomrule
  \end{tabular}
\end{table}

\subsection{Phân đoạn bông lúa đơn}

Đếm số lượng hạt trên mỗi bông lúa thông qua phân tích ảnh và kỹ thuật học sâu vẫn là một
nhiệm vụ đầy thách thức, với nhiều phương pháp vẫn đang được phát triển. Chúng tôi kiểm tra
Mô hình Phân đoạn Bất kỳ (SAM)~\cite{kirillov2023} và Mô hình Phân đoạn Bất kỳ 2
(SAM2)~\cite{ravi2024} để phân đoạn ảnh trong phòng của các bông lúa riêng lẻ. Chúng tôi
chọn dòng mô hình SAM vì phạm vi ứng dụng rộng và không yêu cầu chú thích dữ liệu phức tạp.
Kết quả phân đoạn bởi các mô hình SAM cho thấy, tính trung bình, 83\% hạt lúa trong các bông
được mở ra bằng tay có thể được nhận diện, trong khi chỉ 73\% hạt trong trạng thái treo tự
nhiên có thể được phân đoạn chính xác. Do sự che khuất và chồng lấp giữa các hạt, đặc biệt
ở trạng thái treo tự nhiên, rất khó nhận diện và tách biệt chúng để đếm. Đối với các mô hình
hiện tại, tốt hơn nên mở ra bông lúa bằng tay. Bằng cách cung cấp nhiều ảnh bông lúa ở cả
trạng thái tự nhiên và trạng thái mở ra, cơ sở dữ liệu của chúng tôi cho phép phát triển
chuyên sâu hơn các mô hình AI cho nhiệm vụ kiểu hình thực vật quan trọng nhưng đầy thách
thức này.

\subsection{Tái tạo 3D cây lúa với phân đoạn bông}

Chúng tôi tiếp tục thực hiện tái tạo 3D để chứng minh phạm vi ứng dụng rộng của bộ dữ liệu
và chất lượng cao của các chú thích. Chuỗi ảnh của một cây đơn và các chú thích phân đoạn
của chúng được chọn. Sau đó chúng tôi sử dụng OSTRA~\cite{xu2023}, một framework mã nguồn
mở cho tái tạo mục tiêu 3D và phân đoạn đa cấp, để tái tạo cây với bông được gán nhãn thông
qua các phương pháp stereo đa góc nhìn (MVS). Bằng cách kết hợp chuỗi ảnh và mặt nạ phân
đoạn của chúng, OSTRA thực hiện ước tính tư thế máy ảnh, ước tính độ sâu và ước tính vector
hướng, và tái tạo hoàn toàn cây 3D. Bốn tấm nhựa trong nền rất hữu ích cho việc khớp đặc
trưng và ước tính vị trí. Chúng tôi đã hoàn thành tái tạo 3D cây cho ba giống lúa khác nhau
thông qua ảnh và chú thích của chúng.

% ===================================================================
\section{Thảo luận}

Trong bài báo này, chúng tôi giới thiệu CVRP, một bộ dữ liệu hình ảnh cây lúa đa giống và
đa góc nhìn. CVRP bao gồm 2.303 ảnh đồng ruộng của toàn bộ cây và 123 ảnh bông lúa trong
phòng. Để ghi lại các đặc điểm chi tiết của cây lúa và bông lúa nhằm quan sát sự khác biệt
đặc trưng theo giống, các cây được tách khỏi cây xung quanh và nền đồng trong quá trình chụp
ảnh. Do đó, CVRP không chỉ phục vụ cho mục đích nhận diện và phân đoạn bông dựa trên ảnh
thông thường mà còn cung cấp tài nguyên quý giá cho các nhiệm vụ đầy thách thức như nhận diện
giống lúa thông lượng cao, đếm bông và đếm hạt. Ngoài ra, bộ dữ liệu có thể được sử dụng cho
các nghiên cứu mô hình 3D cây trồng.

Nhiều yếu tố có thể ảnh hưởng đến chất lượng ảnh bông lúa trên đồng ruộng. Bông lúa trải qua
ba giai đoạn sinh trưởng riêng biệt: giai đoạn trỗ bông, giai đoạn nở hoa và giai đoạn chín.
Trong các giai đoạn khác nhau, cây lúa có thể biểu hiện sự khác biệt kiểu hình đáng kể. Một
số giống địa phương thậm chí bị ảnh hưởng nghiêm trọng bởi hiện tượng đổ ngã. Các bông có
thể bị lá che khuất nhiều hơn sau khi hạt chắc vì trọng lượng tăng của bông làm thay đổi
hình dạng cây. Đối với thu thập dữ liệu đồng ruộng về bông lúa, giai đoạn nở hoa, khi tất
cả bông đã trỗ, được khuyến nghị cao. Điều này đảm bảo tính toàn vẹn của đặc điểm bông và
giảm tác động tiêu cực của hiện tượng đổ ngã đến thu thập dữ liệu. Đối với thu thập dữ liệu
đồng ruộng, lựa chọn góc nhìn phụ thuộc vào đặc điểm bông được điều tra. Góc nhìn thẳng đứng
phù hợp nhất để ghi lại mật độ bông. Tuy nhiên, nếu tập trung vào từng cây riêng lẻ, góc
nhìn chếch phù hợp hơn vì nó ghi lại cả số lượng bông lẫn các đặc điểm chi tiết như hình
dạng và kết cấu. Góc nhìn cận cảnh cũng có giá trị để cung cấp các đặc điểm chi tiết của
từng bông riêng lẻ. Nếu không sử dụng thu thập dữ liệu thông lượng cao, chúng tôi khuyến
nghị sử dụng độ phân giải $1080 \times 1920$, đủ cho hầu hết các nghiên cứu.

Để đơn giản hóa quy trình chú thích bông lúa, chúng tôi đã áp dụng quy trình hai giai đoạn
kết hợp nỗ lực thủ công và dự đoán tự động. Trong giai đoạn đầu, một phần ảnh được chú thích
thủ công thông qua công cụ LabelMe. Các ảnh đã chú thích này sau đó được sử dụng làm tập
huấn luyện để huấn luyện mô hình phân đoạn, mô hình này sau đó tạo ra dự đoán cho các ảnh
còn lại. Các chỉnh sửa nhỏ được áp dụng cho các chú thích được dự đoán để đảm bảo độ chính
xác. Phương pháp lai ghép này đã cải thiện đáng kể hiệu quả trong khi duy trì chất lượng chú
thích cao. Khi tạo các bộ dữ liệu liên quan đến công việc chú thích rộng rãi, công cụ dự đoán
chú thích là rất quan trọng. Chúng tôi đánh giá hiệu suất phân đoạn của một số mô hình SOTA
hiện tại trên CVRP và chọn Mask2Former, mô hình hoạt động tốt nhất trong số đó, cho mục đích
của chúng tôi.

Độ chính xác phân đoạn của các mô hình thông thường đối với bông lúa vẫn chưa tối ưu. Sự
tương đồng về hình ảnh giữa bông lúa và lá xung quanh đặt ra một thách thức độc đáo, đòi hỏi
các kỹ thuật trích xuất đặc trưng tinh vi hơn. Để giải quyết những thách thức này, có thể có
lợi khi tinh chỉnh các mô hình lớn như SAM để phù hợp hơn với các nhiệm vụ phân đoạn cụ thể.
Một phương pháp thay thế là thiết kế mô hình tùy chỉnh được điều chỉnh đặc biệt cho nhiệm vụ
phân đoạn bông lúa. Phương pháp này có thể mang lại độ chính xác cao hơn so với các mô hình
phân đoạn đa mục đích. Các phương pháp có mục tiêu như vậy có thể bao gồm việc kết hợp kiến
thức đặc thù theo lĩnh vực và sử dụng các kỹ thuật tiên tiến như cơ chế chú ý hoặc tổng hợp
đặc trưng đa tỷ lệ, nhằm nâng cao hiệu suất.

% ===================================================================
\section*{Tính khả dụng của dữ liệu}

Bộ dữ liệu CVRP được công bố công khai trên Hugging Face tại
\url{https://huggingface.co/datasets/CVRPDataset/CVRP} để sử dụng học thuật theo giấy phép
được quy định. Mã nguồn và các mô hình đã huấn luyện có tại
\url{https://huggingface.co/CVRPDataset/Model}. Bộ dữ liệu, mã nguồn và mô hình đã huấn luyện
cũng có tại \url{http://bic.njau.edu.cn/CVRP.html}.

% ===================================================================
\section*{Đóng góp của tác giả}

X.G. thiết kế nghiên cứu. Z.T. đóng góp vào xử lý dữ liệu, mô hình hóa mạng nơ-ron, chú
thích và kiểm tra ảnh. J.S. đóng góp vào thu thập dữ liệu, chú thích và kiểm tra ảnh. Y.T.
đóng góp vào canh tác cây và thu thập dữ liệu. J.X. đóng góp vào tái tạo mục tiêu 3D. W.Z.
đóng góp vào mô hình hóa và huấn luyện mạng nơ-ron. G.J. đóng góp vào thu thập mẫu cây, và
J.D. đóng góp vào tái tạo mục tiêu 3D. Z.T. và X.G. viết bản thảo. Tất cả tác giả đã đọc và
phê duyệt bản thảo.

% ===================================================================
\section*{Tài trợ}

Công trình này được hỗ trợ một phần bởi các tài trợ từ Dự án Khoa học và Công nghệ Quốc gia
Trọng điểm Nhân giống Sinh học (Mã số 2023ZD04076), Quỹ Khoa học Tự nhiên Quốc gia Trung
Quốc (Mã số 32170647), Quỹ Khoa học Tự nhiên Tỉnh Giang Tô Trung Quốc (Mã số BK20212010 và
BE2022383), Trung tâm Nghiên cứu Kỹ thuật Giang Tô về Chỉnh sửa Genome Thực vật, Công ty
TNHH Nghiên cứu và Phát triển Lúa Japonica Phía Nam, và Trung tâm Đổi mới Hợp tác Giang Tô
về Sản xuất Cây trồng Hiện đại.

% ===================================================================
\section*{Tuyên bố về lợi ích cạnh tranh}

Các tác giả tuyên bố rằng họ không có lợi ích tài chính cạnh tranh hoặc mối quan hệ cá nhân
đã biết nào có thể xuất hiện để ảnh hưởng đến công trình được báo cáo trong bài báo này.

% ===================================================================
\section*{Lời cảm ơn}

Chúng tôi cảm ơn Giáo sư Jianmin Wan vì những đề xuất quý báu và Ông Zhitao Zhu, Ông Yu
Wang, TS. Weijie Tang, và TS. Yunhui Zhang vì sự hỗ trợ kỹ thuật của họ.

% ===================================================================
\begin{thebibliography}{99}

\bibitem{chandra2020}
A.L. Chandra, S.V. Desai, V.N. Balasubramanian, S. Ninomiya, W. Guo,
Học tích cực với giám sát điểm để phát hiện bông lúa hiệu quả về chi phí trong cây ngũ cốc,
\textit{Plant Methods} 16 (2020) 1--16.

\bibitem{xiong2017}
X. Xiong, L. Duan, L. Liu, H. Tu, P. Yang, D. Wu, G. Chen, L. Xiong, W. Yang, Q. Liu,
Panicle-SEG: một phương pháp phân đoạn ảnh mạnh mẽ cho bông lúa ngoài đồng dựa trên học sâu
và tối ưu hóa superpixel,
\textit{Plant Methods} 13 (2017) 1--15.

\bibitem{bai2018}
X. Bai, Z. Cao, L. Zhao, J. Zhang, C. Lv, C. Li, J. Xie,
Quan sát tự động giai đoạn trỗ bông lúa bằng phương pháp phát hiện bông lúa dựa trên tầng
nhiều bộ phân loại,
\textit{Agric. For. Meteorol.} 259 (2018) 260--270.

\bibitem{hayat2020}
M.A. Hayat, J. Wu, Y. Cao,
Học Bayesian không giám sát để phân đoạn bông lúa với ảnh UAV,
\textit{Plant Methods} 16 (1) (2020) 18.

\bibitem{xiao2023}
L. Xiao, Z. Pan, X. Du, W. Chen, W. Qu, Y. Bai, T. Xu,
Kết hợp đặc trưng skip-connection có trọng số: một phương pháp tăng cường phân đoạn ảnh bông
lúa theo định hướng UAV,
\textit{Comput. Electron. Agric.} 207 (2023) 107754.

\bibitem{zhou2023}
Q. Zhou, W. Guo, N. Chen, Z. Wang, G. Li, Y. Ding, S. Ninomiya, Y. Mu,
Phân tích tác động của nitơ đến sự phát triển bông lúa bằng phát hiện bông và theo dõi theo
chuỗi thời gian,
\textit{Plant Phenomics} 5 (2023) 48.

\bibitem{wang2021}
H. Wang, S. Lyu, Y. Ren,
Bộ dữ liệu ảnh lúa nước để phân đoạn bông,
\textit{Agronomy} 11 (8) (2021) 1542.

\bibitem{david2020}
E. David và cộng sự,
Bộ dữ liệu phát hiện đầu lúa mì toàn cầu (GWHD): một bộ dữ liệu ảnh RGB độ phân giải cao
được gán nhãn đa dạng và chất lượng cao để phát triển và đánh giá chuẩn các phương pháp phát
hiện đầu lúa mì,
\textit{Plant Phenomics} 2020 (2020) 12.

\bibitem{david2021}
E. David và cộng sự,
Phát hiện đầu lúa mì toàn cầu 2021: bộ dữ liệu cải tiến để đánh giá chuẩn các phương pháp
phát hiện đầu lúa mì,
\textit{Plant Phenomics} 2021 (2021) 9.

\bibitem{hani2020}
N. H{\"a}ni, P. Roy, V. Isler,
MinneApple: một bộ dữ liệu chuẩn để phát hiện và phân đoạn táo,
\textit{IEEE Rob. Autom. Lett.} 5 (2) (2020) 852--858.

\bibitem{vgg2016}
A. Dutta, A. Gupta, A. Zissermann,
VGG image annotator (VIA),
[Online]. Có tại: \url{http://www.robots.ox.ac.uk/vgg/software/via/}, 2016.

\bibitem{russell2008}
B.C. Russell, A. Torralba, K.P. Murphy, W.T. Freeman,
LabelMe: cơ sở dữ liệu và công cụ chú thích ảnh dựa trên web,
\textit{Int. J. Comput. Vis.} 77 (2008) 157--173.

\bibitem{mmseg2020}
Mms Contributors,
MMSegmentation: bộ công cụ và chuẩn phân đoạn ngữ nghĩa OpenMMLab,
[Online]. Có tại: \url{https://github.com/open-mmlab/mmsegmentation}, 2020.

\bibitem{kirillov2023}
A. Kirillov và cộng sự,
Segment anything,
trong: \textit{Proceedings of the IEEE/CVF International Conference on Computer Vision},
2023, tr. 4015--4026.

\bibitem{ravi2024}
N. Ravi và cộng sự,
SAM 2: Segment Anything in Images and Videos.

\bibitem{xu2023}
J. Xu, W. Zhao, Z. Tang, X. Gan,
Phương pháp tái tạo mục tiêu 3D và phân đoạn đa cấp một điểm dừng,
\textit{ArXiv Prepr.} ArXiv2308.06974 (2023).

\bibitem{chen2018}
L.-C. Chen, Y. Zhu, G. Papandreou, F. Schroff, H. Adam,
Bộ mã hóa-giải mã với tích chập phân tách atrous để phân đoạn ảnh ngữ nghĩa,
trong: \textit{Proceedings of the European Conference on Computer Vision (ECCV)},
2018, tr. 801--818.

\bibitem{cheng2022}
B. Cheng, I. Misra, A.G. Schwing, A. Kirillov, R. Girdhar,
Masked-attention mask transformer để phân đoạn ảnh phổ quát,
trong: \textit{Proceedings of the IEEE/CVF Conference on Computer Vision and Pattern Recognition},
2022, tr. 1290--1299.

\bibitem{xie2021}
E. Xie, W. Wang, Z. Yu, A. Anandkumar, J.M. Alvarez, P. Luo,
SegFormer: thiết kế đơn giản và hiệu quả cho phân đoạn ngữ nghĩa với transformer,
\textit{Adv. Neural Inf. Process. Syst.} 34 (2021) 12077--12090.

\bibitem{zhang2021}
W. Zhang, J. Pang, K. Chen, C.C. Loy,
K-net: hướng tới phân đoạn ảnh thống nhất,
\textit{Adv. Neural Inf. Process. Syst.} 34 (2021) 10326--10338.

\end{thebibliography}

\end{document}
